% page 96 du fac-simile (image p-095 du scan)
% bloc 154.9 x 233.3 mm ; marge gauche 8.0 mm
\pgc{-12.700mm}{0.000mm}{
\l{\cel{6}88}
\l{}
\l{\sou{cherpar}.\cel{2}(\sou{trans., ek, de}). Prenar ek puteo, fonto, barelo,}
\l{\cel{3}e c. kelke de la liquido quan lu kontenas. - (\sou{anke metaf.})}
\l{}
\l{}
\l{\sou{chevrono}. I. Peco ek ligno, fixigita an la pento di tekto,}
\l{\cel{3}e suportanta la lati qui subtenas la tekto-kovruro, la}
\l{\cel{3}teguli, la ardezo-plaki. - II. Bendi plata, dispozita}
\l{\cel{3}V-forme sur la skudo, la blazoni. - EFRS.}
\l{}
\l{\sou{chifrar}. (\sou{trans.)}\cel{2}Kompozar kripto-grafuro, t.e. texto sekre-}
\l{\cel{3}ta, qua tradukas la vera texto per ula moyeno o regulo}
\l{\cel{3}(\textquotedbl{}klefo\textquotedbl{}). - DEF.}
\l{}
\l{\sou{chigo}. (\sou{zool.}) Insekto di Sud-Amerika. - L.\cel{1}\sou{sarcopsylla pe}-}
\l{\cel{3}\sou{netrans}). - EFS.}
\l{}
\l{\sou{chika}. Qua havas gracio eleganta, audacoza. - DEFR.}
\l{}
\l{\sou{chilo}. (\sou{Fiziol}.)Suko qua konsistas (en la intestino tenua)}
\l{\cel{3}ye la parto nutriva di la bulo nutrala, qua divenas sango.}
\l{\cel{3}- DEFIRS.}
\l{}
\l{\sou{chimo}. Sorto di paplo quan produktas, en la stomako, la elu-}
\l{\cel{3}boro komencala di la nutrivi. - DEFIRS.}
\l{}
\l{\sou{chimpanzeo}. (\sou{zool}.) Simio alta-statura. L.\cel{1}\sou{troglodytea niger}.}
\l{\cel{3}- DEFIRS.}
\l{}
\l{\sou{chiniono}.. Che la muliero : la parto di la hararo qua amas-}
\l{\cel{3}igesas ed erektesas dop la kapo. - DEF.}
\l{}
\l{\sou{chipa}. Qua kustas min multe kam lo analoga. - E.}
\l{}
\l{\sou{chitino}. (\sou{kemio}).Elemento precipua di la kuraso di la \textquotedbl{}artro-}
\l{\cel{3}pedi\textquotedbl{}. -\cel{1}\sou{Simbolo kemiala}\cel{1}: C16, H26, N2, O10. - DEFIS.}
\l{}
\l{\sou{chitozamino}. (\sou{kemio)}. Bazo qua formacesas per hidrolizo di}
\l{\cel{3}la chitino e di multa albuminoidi.}
\l{}
\l{\sou{choano}. (\sou{anat.}) Orifico dopa di la naztrui. - DEFIS.}
\l{}
\l{\sou{chokolado}. Substanco nutrala, ek kakao-grani, pistita kun}
\l{\cel{3}sukro. - DEFIRS.}
\l{}
\l{\sou{chomar}. (\sou{netrans.}) Nelaborar pro neemployateso nevolata. - F}
\l{}
\l{\sou{chopino}. Birglaso qua kontenas cirkume un mi-litro. - DEF.}
\l{}
\l{\sou{choseo}. I. Bendo de tereno, ofte stonizita, qua dominacas}
\l{\cel{3}rivero, lageto, marsho, quan lu transiras, ed uzata kom}
\l{\cel{3}voyo, paseyo. -II. La mezo (generale pavizita o stoniz-}
\l{\cel{3}ita) di voyo, strado. - DF.}
}
